\section{A fully replenished physical obstruction}\label{sec:construction}
Two mandatory depot-to-depot services consume 15 kWh each: A occupies
$[0,1]$ and B occupies $[2,3]$. At most two identical buses are available,
each with 20 kWh capacity, zero reserve and full initial inventory. Early
charging is available only in $[1,2]$ at 10 kW, both per bus and in total.
Terminal charging is available only in $[3,4]$, using one connector of
power $K\le30$ kW. Every used bus must finish full at hour 4. Charging is
continuous and lossless, with no taper or switching time. Travel cost,
movement consumption and auxiliary load are zero. All prices and costs in
this section are synthetic.

There are exactly two service partitions up to vehicle labels. With one
bus, inventory after A is 5; it needs at least 10 early kWh to complete B.
The early limit therefore fixes load $(10,20)$. With two buses, A's bus
buys $x\in[0,10]$ early; terminal purchases are $15-x$ for A's bus and 15
for B's bus. Each bus returns full. Both partitions buy exactly 30 kWh;
initial inventory cannot subsidize the horizon.

Let $f\ge0$ be the cost of a used bus and define
\begin{equation}\label{eq:supply}
 F(x,z)=4x+\frac{x^2+z^2}{10},\qquad x,z\ge0,
 \qquad g(x)=F(x,30-x)=90-2x+\frac{x^2}{5}.
\end{equation}
The early/late distinction permits different marginal supply costs despite
constant total energy.

\begin{theorem}[Nominal failure]\label{thm:nominal}
For $K=30$ and $f=7$,
\[
 D=97,\qquad \CH=\frac{7591}{80}=94.8875,\qquad
 \Delta=\frac{169}{80}>0.
\]
The unique physical planner load is $(10,20)$ with one bus. Its own-price
regret at $p=(6,4)$ is 13.
\end{theorem}

The one-bus objective is $7+g(10)=97$; the two-bus branch minimizes
$14+g(x)$ at $x=5$, giving 99. The hull instead mixes one bus at early
energy 10 with two buses at early energy 0. Its cheapest operating cost at
mean early energy $x$ is $14-0.7x$. Thus its objective is
\[
 14-0.7x+g(x)=\frac{7591}{80}+\frac15(x-27/4)^2.
\]
The mixture weights are $27/40$ and $13/40$. Both components are physical;
the mean operating cost $371/40$ is a convexification. Running those plans
on different days would have average physical cost 99.275, not 94.8875.

At the physical optimum's marginal price $(6,4)$, its bill plus operating
cost is 147. The two-bus zero-early alternative costs 134 at that fixed
tariff. This 13-unit incentive changes the entire fleet. By contrast, the
hull price $(107/20,93/20)$ makes the two hull components tie, but is not
the marginal supply price of the one-bus physical load. The distinction
explains why a hull price and an own-load tariff answer different questions.

\subsection{An exact connector and fee characterization}
Define $\clamp(v;\ell,u)=\min\{u,\max\{\ell,v\}\}$. The next result
describes the entire declared parameter domain, including its boundaries.

\begin{theorem}[Connector capacity and vehicle fee]\label{thm:capacity}
For $0\le K\le30$ and $f\ge0$, the model is infeasible if $K<20$.
If $K=20$, $D=\CH=f+90$; both service partitions tie only at $f=0$.
For $K>20$, put
\[
 \ell=30-K,\quad d=K-20,\quad t=\max\{5,\ell\},\quad
 h=\clamp\left(5+\frac{5f}{2d};\ell,10\right).
\]
Then
\begin{align}
 D&=\min\{f+90,\ 2f+g(t)\},\label{eq:capacityD}\\
 \CH&=2f-\frac{f(h-\ell)}d+g(h).\label{eq:capacityCH}
\end{align}
For $f>0$, the gap is strictly positive exactly when
\begin{equation}\label{eq:capacityregion}
 \max\left\{0,\frac25d(\ell-5)\right\}<f<2d.
\end{equation}
For $f=0$ the gap is zero. Both physical planner partitions are retained
when $f=90-g(t)$.
\end{theorem}

The proof is a physical projection and one-dimensional convex minimization,
not a finite parameter scan (Appendix~\ref{app:capacity}). In particular,
with $f=7$ the gap is zero for $20\le K\le47/2$ and positive for
$47/2<K\le30$. Extra charging capacity can create a price-support failure
by enlarging the alternative set; it does not increase the planner's minimum
cost or show that installing a connector worsens welfare.

The full price-response value is
\begin{equation}\label{eq:fullresponse}
 V(p)=\min\{f+10p_1+20p_2,\;
 2f+\ell p_1+(30-\ell)p_2,\;2f+10p_1+20p_2\}.
\end{equation}
An affine objective attains its minimum on each physical segment at an
endpoint. If $p_1=p_2$, its whole two-bus interval is minimizing within
that partition and must be retained whenever its bill also ties globally.
This includes the zero-fee boundaries; an endpoint representative alone
would omit part of the response set.

\subsection{Positive reserve and charging losses}
The nominal early limit is binding: adding reserve or losses without
changing power can remove the one-bus branch. This is a feasibility issue,
not disappearance of the pricing mechanism. A nearby construction uses
reserve 1 kWh, efficiency $19/20$, early power 12 kW and one terminal
connector of 30 kW, keeping the service energies, capacity 20 and fee 7.
The one-bus state sequence is
\[
 20\longrightarrow5\longrightarrow16\longrightarrow1
 \longrightarrow20.
\]
It buys $220/19$ grid kWh early and 20 terminally. A physical two-bus
endpoint buys $30/19$ early and 30 terminally; terminal sessions of
$9/19$ and $10/19$ hour fit serially.

\begin{proposition}[A nondegenerate neighborhood]\label{prop:reserve}
For these modified parameters,
\[
 D=\frac{38527}{361},\qquad
 \CH=\frac{2987911}{28880},\qquad
 \Delta=\frac{94249}{28880}>0.
\]
With terminal power 30 fixed, a positive gap persists in an open
neighborhood of the stated reserve, efficiency and early power.
\end{proposition}
Appendix~\ref{app:reserve} states the strict inequalities defining this
neighborhood. The result remains conditional on constant efficiency,
zero switching time and the absence of charging taper.
